\BourbakiExerciseGroup

\begin{enumerate}
\def\labelenumi{\arabic{enumi})}
\item
  Show that the countable Hausdorff space defined in § 9, Exercise 20 c), is connected.
\item
  \begin{enumerate}
  \def\labelenumii{\alph{enumii})}
  \tightlist
  \item
    Let X be a topological space and \(\mathfrak{C}\) the topology of X. Show that if X is connected in the semi-regular topology \(\mathfrak{C}^*\) associated with \(\mathfrak{C}\) ( \(\S 8\) , Exercise 20), then X is connected in the original topology \(\mathfrak{C}\) . Deduce that there exist connected submaximal spaces ( \(\S 8\) , Exercise 22).
  \end{enumerate}
\end{enumerate}

*b) In particular construct a topology on the real line \(\mathbf{R}\) , with respect to which \(\mathbf{R}\) is connected but no point of \(\mathbf{R}\) has a fundamental system of connected neighbourhoods {[}cf.~Chapter IV, § 2, Exercise 14 c){]}.*

\begin{enumerate}
\def\labelenumi{\arabic{enumi})}
\setcounter{enumi}{2}
\tightlist
\item
  Let A, B be two subsets of a topological space X.
\end{enumerate}

\begin{enumerate}
\def\labelenumi{\alph{enumi})}
\item
  Show that if A and B are closed in X, and A ∪ B and A ∩ B are connected, then A and B are connected. Show by an example that this statement is no longer necessarily true when one of the sets A, B is not closed.
\item
  Suppose that A and B are connected and \(\overline{A} \cap B \neq \varnothing\) . Show that \(A \cup B\) is connected.
\end{enumerate}

\begin{enumerate}
\def\labelenumi{\arabic{enumi})}
\setcounter{enumi}{3}
\tightlist
\item
  Let X be a connected space of at least two points.
\end{enumerate}

\begin{enumerate}
\def\labelenumi{\alph{enumi})}
\item
  Let A be a connected subset of X, B a subset of \(\mathbb{C}\mathbf{A}\) which is both open and closed in \(\mathbb{C}\mathbf{A}\) . Show that \(\mathbf{A} \cup \mathbf{B}\) is connected. (Apply Exercise 5 of §3 to \(\mathbf{Y} = \mathbf{A} \cup \mathbf{B}\) and \(\mathbf{Z} = \mathbb{C}\mathbf{A}\) .)
\item
  Let A be a connected subset of X and B a component of the set CA. Show that CB is connected {[}use a){]}.
\item
  Deduce from b) that there exist two connected subsets M, N of X which are distinct from X and such that \(M \cup N = X\) and \(M \cap N = \varnothing\) .
\end{enumerate}

*5) a) Give an example, in the real plane \(\mathbf{R}^2\) , of a decreasing sequence \((\mathbf{A}_n)\) of connected sets whose intersection is not connected (cf.~Chapter II, § 4, Exercise 14).

\begin{enumerate}
\def\labelenumi{\alph{enumi})}
\setcounter{enumi}{1}
\tightlist
\item
  In \(\mathbf{R}^2\) , let \(\mathbf{X}\) be the set which is the union of the open half-planes \(y > 0\) and \(y < 0\) and the points \((n, 0)\) where \(n\) is an integer \(> 0\) .
\end{enumerate}

X is connected with respect to the topology \(\mathfrak{C}_0\) induced by that of \(\mathbb{R}^2\) . Define a sequence \((\mathfrak{C}_n)\) of topologies on X such that \(\mathfrak{C}_{n+1}\) is finer than \(\mathfrak{C}_n\) , and such that X is connected with respect to each \(\mathfrak{C}_n\) but not connected with respect to the topology which is the least upper bound of the sequence \((\mathfrak{C}_n)\) {[}pass from \(\mathfrak{C}_n\) to \(\mathfrak{C}_{n+1}\) by introducing new neighbourhoods of the point \((n+1,0)]\) .*

\begin{enumerate}
\def\labelenumi{\arabic{enumi})}
\setcounter{enumi}{5}
\item
  Let X, Y be two connected spaces and let A (resp. B) be a subset of X (resp. Y) distinct from X (resp. Y). Show that in the product space \(X \times Y\) the complement of \(A \times B\) is connected.
\item
  Let X, Y be two connected spaces and f a mapping of \(X \times Y\) into a topological space Z such that all the partial mappings \(f(0, y) : X \to Z\) , \(f(x, 0) : Y \to Z (x \in Y, y \in Y)\) are continuous. Show that \(f(X \times Y)\) is connected.
\item
  In the example of a topology on a product set \(\mathbf{X} = \prod_{i\in I}\mathbf{X}_i\) given in § 4, Exercise 9, suppose that I is infinite and the cardinal t uncountable. Show that if the \(\mathbf{X}_i\) are connected, regular, and have each at least two distinct points, then the connected component of a point \(a = (a_i)\) of X is the set of \(x = (x_i)\) such that \(x_{i} = a_{i}\) except for a finite number of indices \(i\in I\) . (Reduce to the case I = N; consider two points \(b = (b_n)\) , \(c = (c_n)\) of X such that \(b_{n}\neq c_{n}\) for all n.~For each n let \((\mathrm{V}_{nm})_{m\geqslant 0}\) be a sequence of open neighbourhoods of \(b_{n}\) in \(\mathbf{X}_n\) which do not contain \(c_{n}\) and are such that \(\overline{\mathbf{V}}_{n,m + 1}\subset \mathbf{V}_{nm}\) . Let Z be the set of points \((x_{n})\) of X with the following property: there exists an integer \(k > 0\) such that \(x_{n}\in \mathrm{V}_{n,n - k}\) whenever \(n\geqslant k\) . Show that Z is both open and closed in X).
\end{enumerate}

*9) Show that in the locally compact space \(\mathbf{X}\) defined in § 10, Exercise 20, there are points \(x\) whose component in \(\mathbf{X}\) is different from the intersection of the subsets of \(\mathbf{X}\) which are both open and closed and contain \(x_*\)

\begin{enumerate}
\def\labelenumi{\arabic{enumi})}
\setcounter{enumi}{9}
\item
  Show that a linearly ordered set X, with the topology \(\mathcal{G}_{+}(X)\) or the topology \(\mathcal{G}_{-}(X)\) (§ 2, Exercise 5) is totally disconnected.
\item
  Let X be a topological space and R an equivalence relation on X such that every equivalence class mod R is contained in a component of X. Show that the components of X/R are the canonical images of the components of X (cf.~Chapter III, § 2, Exercise 17).
\item
  Let X be a topological space. Show that the following three conditions are equivalent: \(\alpha\) ) the components of X are open sets; \(\beta\) ) every \(x \in X\) has a neighbourhood contained in every set which is both open and closed and contains \(x; \gamma\) ) for every \(x \in X\) , the intersection of the sets which are both open and closed and contain x is an open set.
\end{enumerate}

*13) Let \(\theta\) be an irrational number. Let \(f\) be the mapping of the real line \(\mathbf{R}\) into the torus \(\mathbf{T}^2\) for which the image \(f(t)\) of each \(t \in \mathbf{R}\) is the canonical image in \(\mathbf{T}^2\) of \((t, \theta t) \in \mathbf{R}^2\) . \(f\) is injective and continuous. Show that, in the topology on \(\mathbf{R}\) which is the inverse image under \(f\) of the topology of \(\mathbf{T}^2\) , \(\mathbf{R}\) is connected but no point of \(\mathbf{R}\) has a fundamental system of connected neighbourhoods.*

\begin{enumerate}
\def\labelenumi{\arabic{enumi})}
\setcounter{enumi}{13}
\tightlist
\item
  \begin{enumerate}
  \def\labelenumii{\alph{enumii})}
  \tightlist
  \item
    Let X be a locally compact and locally connected space, and let A be a closed subset of X. Let A* be the union of A and the components of CA which are relatively compact. A is said to be full if A* = A. Show that A* is closed and full.
  \end{enumerate}
\end{enumerate}

\begin{enumerate}
\def\labelenumi{\alph{enumi})}
\setcounter{enumi}{1}
\item
  Suppose in addition that X is connected. Show that if A is a compact subset of X, then A* is compact (consider a compact neighbourhood V of A and the components of CA which meet the frontier of V).
\item
  Suppose in addition that X is connected and \(\sigma\) -compact. Show that there exists an increasing sequence \((\mathbf{K}_{n})_{n\geqslant0}\) of full connected compact sets whose union is X and which are such that \(K_{n}\subset\mathring{K}_{n+1}\) for all n.
\end{enumerate}

*d) Show that the conclusion of b) is no longer valid if we suppose only that X is connected and locally compact, or that X is Hausdorff, connected and locally connected. {[}Consider the subspace of \(\mathbf{R}^2\) consisting of points \((x, y)\) such that either \(y = 0\) and \(0 \leqslant x \leqslant 1\) , or else \(y > 0\) and \(x = 1/n\) where \(n\) is an integer \(> 0\) .*

¶ 15) Let X be a connected, locally connected space.

\begin{enumerate}
\def\labelenumi{\alph{enumi})}
\item
  Let M, N be two disjoint non-empty closed subsets of X. Show that \(\mathbb{C}(M \cup N)\) has a component whose closure meets both M and N. {[}If C is a component of \(\mathbb{C}(M \cup N)\) whose closure does not meet M, observe that C is both open and closed in CN{]}.
\item
  If A is a closed subset of X distinct from X, deduce from a) that every component of A meets the closure of \(\mathbb{C}\mathbb{A}\) .
\end{enumerate}

*16) Let X be the subspace of the real plane \(\mathbf{R}^2\) which is the union of the line \(y = 0\) , the closed segments with end points (0, 1) and \((n, 1/(1 + n))\) , and the half-lines \(x \leqslant n\) , \(y = 1/(1 + n)\) ( \(n\) an integer \(>0\) ). Show that X is a connected space but that there is a component C of the complement of (0, 1) in X such that (0, 1) \(\notin \overline{\mathbf{C}}\) {[}(compare with Exercises 4 a) and 15 b){]}.*

\(^{*}\) 17) Let \(I_{0}\) be the interval \([-1,1]\) of R and let \(P_{0}\) be the quotient space of \(I_{0}\) obtained by identifying the two points 1/2 and 1; let \(B_{0}\) be the quotient space of \(I_{0}\) obtained by identifying the points 1/2 and 1 and also identifying -1/2 and -1. Let I be the open interval {]}--- I, I{[} in R, and let P (resp. B) be the canonical image of I in \(P_{0}\) (resp. \(B_{0}\) ).

\begin{enumerate}
\def\labelenumi{\alph{enumi})}
\item
  Show that no two of the three spaces I, P, B are homeomorphic (to show that I and P are not homeomorphic, notice that there are points in P whose complement is connected; similarly for the other pairs of spaces).
\item
  Let X be the sum of a countably infinite family of spaces homeomorphic to I and a countably infinite family of spaces homeomorphic to B. Let Y be the sum of X and P. Show that X and Y are not homeomorphic, but that (i) there is a continuous bijection of X onto Y and a continuous bijection of Y onto X; (ii) X is homeomorphic to an open subspace of Y and Y is homeomorphic to an open subspace of X.*
\end{enumerate}

*18) Let X be the real plane \(\mathbf{R}^2\) , and let R be the equivalence relation on X whose equivalence classes are the lines \(y = \beta\) for \(\beta > 0\) and the half-lines \(x = \alpha\) , \(y \leqslant 0\) for all \(\alpha \in \mathbb{R}\) . Show that if X/R is regarded as a subset of \(\mathfrak{P}_0(X)\) , then the topology induced on X/R by \(\mathcal{C}_{\Omega}\) ( \(\S 2\) , Exercise 7) is discrete, and that the topology induced on X/R by \(\mathcal{C}_{\Phi}\) is not discrete but makes X/R into a non-connected space. Show finally that X/R is connected in the quotient topology.*

*¶ 19) Let X be a locally compact space. It can be shown that X can be identified with a dense open subspace of the universal compact space \(\tilde{X}\) constructed in § 9, Exercise 25 (cf.~Chapter IX, § 1, Exercise 8). Suppose also that X is connected and locally connected.

\begin{enumerate}
\def\labelenumi{\alph{enumi})}
\item
  Show that, for every compact subset K of X, the set of components of X-K which are not relatively compact is finite {[}cf.~Exercise 14 b){]}. Deduce that a point of \(\tilde{X} - X\) belongs to the closure of at most one of the components of X-K {[}if \(A_i\) ( \(i \leqslant i \leqslant n\) ) are the non-compact components of X-K, and \(C_i\) the compact set of points of \(\tilde{X} - X\) which lie in the closure of \(A_i\) , construct a compact space whose underlying set is the sum of X and the sets \(C_i\) , such that X is identified with a dense open subspace of this space{]}.
\item
  Let \(R'\{x, y\}\) be the following equivalence relation on \(\tilde{X}-X\) : ``for every compact subset K of X, x and y lie in the closure of the same non-compact component of X-K''. Let R be the equivalence relation on \(\tilde{X}\) whose equivalence classes are the classes mod \(R'\) and the subsets consisting of a single point of X. Show that R is Hausdorff and that X may be identified with a dense open subspace of the compact space \(X^{b}=\tilde{X}/R\) ; show that in the compact space \(X^{b}-X\) every point is the intersection of its neighbourhoods which are both open and closed. The points of \(X^{b}-X\) are called the ends of the space X.
\item
  Let Y be a compact space such that X may be identified with a dense open subspace of Y and such that in the compact space Y-X every point is the intersection of its neighbourhoods which are both open and closed (*). Show that there is a continuous mapping of \(X^{b}\) into Y which induces the identity mapping on X. {[}If Y-X is the union of two disjoint non-empty closed sets A and B, show that there exist two disjoint open subsets U, Y in V such that \(A \subset U\) , \(B \subset V\) and such that \(Y - (U \cup V)\) is compact{]}.
\item
  Show that if X is σ-compact, then \(X^{b}-X\) has a countable base.
\item
  The real line has two ends, and \(R^{n} (n \geqslant 2)\) has one. Give an example of a connected, locally connected, locally compact subspace of \(R^{2}\) for which the set of ends has the power of the continuum.*
\end{enumerate}

\begin{enumerate}
\def\labelenumi{\arabic{enumi})}
\setcounter{enumi}{19}
\tightlist
\item
  Let X be a Hausdorff space in which each point has a fundamental system of neighbourhoods which are both open and closed in X. Let Y be a subspace of X, and let A be a compact subset of Y which is both open and closed in Y; show that there is a subset B of X which is both open and closed in X, such that \(B \cap Y = A\) .
\end{enumerate}

¶ 21) a) Show that the following conditions on a topological space X are equivalent: (i) for each open set U in X, \(\overline{U}\) is open; (ii) for each closed set F in X, \(\dot{F}\) is closed; (iii) for each disjoint pair of open sets U, V in X, we have \(\overline{U} \cap \overline{V} = \varnothing\) . A Hausdorff space which satisfies these conditions is said to be extremely disconnected; it is then totally disconnected. *The rational line is totally disconnected but not extremely disconnected.*

\begin{enumerate}
\def\labelenumi{\alph{enumi})}
\setcounter{enumi}{1}
\item
  A Hausdorff space is extremely disconnected if and only if the associated semi-regular space (§ 8, Exercise 20) is extremely disconnected.
\item
  Let X be a Hausdorff space and let A be a dense subset of X. Show that if A is extremely disconnected and if the topology of X is A-maximal (§ 3, Exercise 11), then X is extremely disconnected.
\item
  Deduce from b) and c) that, if X is any infinite set, the (compact) ultrafilter space of X is extremely disconnected (cf.~§ 9, Exercises 26 and 24) (**).
\item
  Let X be a Hausdorff space with no isolated points. Show that the topology of X is quasi-maximal (§ 2, Exercise 6) if and only if X is submaximal (§ 8, Exercise 22) and extremely disconnected. (To show that the condition is sufficient, use Exercise 22 e) of § 8 and prove that every closed subset of X which has no isolated points is open).
\item
  Show that every ultraregular space (§ 8, Exercise 25) is extremely disconnected, but that an extremely disconnected compact space with no isolated points is not ultraregular (cf.~§ 9, Exercise 27).
\item
  Show that an extremely disconnected semi-regular space (§ 8, Exercise 20) is regular (cf.~Chapter II, § 4, Exercise 12).
\item
  Show that in an extremely disconnected space there is no convergent sequence which has an infinite number of distinct terms {[}argue by contradiction by constructing recursively two sequences \((\mathbf{U}_{n})\) , \((\mathbf{V}_{n})\) of mutually disjoint open sets such that if \(U = \bigcup_{n} U_{n}\) and \(V = \bigcup_{n} V_{n}\) then \(U \cap V = \varnothing\) but \(\overline{U} \cap \overline{V} \neq \varnothing\) {]}.
\item
  Show that in an extremely disconnected space X a non-isolated point a cannot have a fundamental system of simultaneously open and closed neighbourhoods which is well-ordered (with respect to the relation \(\supset\) ). (Use reductio ad absurdum, as in h).
\end{enumerate}

¶ 22) a) Show that in an extremely disconnected space every open subspace and every dense subspace is extremely disconnected.

\begin{enumerate}
\def\labelenumi{\alph{enumi})}
\setcounter{enumi}{1}
\item
  Let X be a Hausdorff space in which every point has a neighbourhood which is an extremely disconnected subspace of X. Show that X is extremely disconnected.
\item
  Give an example of a Hausdorff space X which is not extremely disconnected but which has an extremely disconnected dense open subspace (paste together two extremely disconnected spaces).
\item
  Let X be a countable discrete space, \(\tilde{X}\) the (compact) space of ultrafilters on X (§ 9, Exercise 26). Let \((\mathbf{A}_{n})_{n\in\mathbb{N}}\) be a countably infinite partition of X into infinite subsets; in the closed subspace \(Y=\tilde{X}-X\) of \(\tilde{X}\) , the sets \(B_{n}=\tilde{X}_{n}\cap Y\) are open and closed and pairwise disjoint. Let \(B=\bigcup_{n}B_{n}\) ; show that the closure B of B in Y is not open in Y, and hence that Y is not extremely disconnected (*). (Use reductio ad absurdum: using Exercise 20, let C be an open and closed subset of \(\tilde{X}\) such that \(C\cap Y=B\) ; consider a point \(x_{n}\) in each of the sets \(C\cap A_{n}\) , let \(\bar{J}\) be the closure of the set J of points \(x_{n}\) , and show that \(\bar{J}\) does not meet B).
\end{enumerate}

\begin{enumerate}
\def\labelenumi{\arabic{enumi})}
\setcounter{enumi}{22}
\item
  Give an example of a bijective mapping \(f\) of a Hausdorff space \(\mathbf{X}\) onto a Hausdorff space \(\mathbf{Y}\) , such that the image under \(f\) of every compact subset of \(\mathbf{X}\) is a compact subset of \(\mathbf{Y}\) , and the image under \(f\) of every connected subset of \(\mathbf{X}\) is a connected subset of \(\mathbf{Y}\) , but such that \(f\) is not continuous (cf.~§ 9, Exercise 4).
\item
  Let X be a topological space.
\end{enumerate}

\begin{enumerate}
\def\labelenumi{\alph{enumi})}
\tightlist
\item
  Let the set \(\mathfrak{P}_{\Omega}(\mathbf{X})\) carry one of the topologies \(\mathfrak{C}_{\Omega}, \mathfrak{C}_{\Phi}\) ( \(\S 2\) , Exercise 7) or \(\mathfrak{C}_{\Theta}\) ( \(\S 8\) , Exercise 12). Show that if a subset \(\mathfrak{B}\) of \(\mathfrak{P}_{\Omega}(\mathbf{X})\) is connec-
\end{enumerate}

ted and if each \(\mathbf{M} \in \mathfrak{B}\) is connected, then \(\bigcup_{\mathbf{M} \in \mathfrak{M}} \mathbf{M}\) is connected.

\begin{enumerate}
\def\labelenumi{\alph{enumi})}
\setcounter{enumi}{1}
\tightlist
\item
  Let \(\mathfrak{S}\) be a subset of \(\mathfrak{P}_{0}(\mathbf{X})\) which contains the set of all finite non-empty subsets of \(\mathbf{X}\) . Show that if \(\mathbf{X}\) is connected then so is \(\mathfrak{S}\) {[}use Exercise 7 b) of § 2 and Proposition 8 of § 4, considering the mappings \((x_{1},\ldots ,x_{n})\to \{x_{1},\ldots ,x_{n}\} ]\) .
\end{enumerate}

\begin{enumerate}
\def\labelenumi{\arabic{enumi})}
\setcounter{enumi}{24}
\tightlist
\item
  Given two topological spaces X, Y, a mapping \(f: \mathbf{X} \to \mathbf{Y}\) is said to be a local homeomorphism if for each \(x \in \mathbf{X}\) there is a neighbourhood U of \(x\) such that \(f(\mathbf{U})\) is a neighbourhood of \(f(x)\) in Y, and the mapping \(\mathbf{U} \to f(\mathbf{U})\) which coincides with \(f\) on U is a homeomorphism. Every local homeomorphism is a continuous open mapping.
\end{enumerate}

\begin{enumerate}
\def\labelenumi{\alph{enumi})}
\item
  Suppose that X is Hausdorff and Y locally connected. Let \(f: \mathbf{X} \to \mathbf{Y}\) be a local homeomorphism with the property that there is an integer \(n > 0\) such that \(\overline{f}^1(y)\) has exactly \(n\) points for every \(y \in \mathbf{Y}\) . Show that every \(y \in \mathbf{Y}\) has an open neighbourhood V such that \(\overline{f}^1(V)\) has \(n\) components \(\mathbf{U}_i\) ( \(1 \leqslant i \leqslant n\) ) and such that the mapping \(\mathbf{U}_i \to \mathbf{V}\) which coincides with \(f\) on \(\mathbf{U}_i\) is a homeomorphism. \(f\) is then a proper mapping.
\item
  Suppose that X is Hausdorff and that Y is connected and locally connected. Let \(f: \mathbf{X} \to \mathbf{Y}\) be a proper local homeomorphism. Show that \(f\) satisfies the hypotheses of \(a\) ). {[}If for each integer \(n > 0\) , \(\mathbf{Y}_n\) is the set of \(y \in \mathbf{Y}\) such that \(\bar{f}^1(y)\) has at least \(n\) points, show that \(\mathbf{Y}_n\) is both open and closed in \(\mathbf{Y}\) {]}.
\end{enumerate}

\(^{*}c)\) Give an example of a surjective local homeomorphism \(f: X \to Y\) , in which Y is compact, connected and locally connected, X is locally compact, connected and locally connected, \(\overline{f}^{1}(y)\) consists of at most two points for each \(y \in Y\) , but f is not proper (cf.~§ 5, Exercise 4).
% semantic-fragments: legacy-input-seam
\relax{}
